\begin{table}[t]
\centering
\caption{Cross-platform validation on Gymnasium-Robotics FetchReach-v4. Rule-based policies with the same FSM backbone and memory variants as Table~\ref{tab:main_results}, evaluated under FO and WeakPO. Mean over 3 seeds $\times$ 60 episodes. The same layered pattern emerges: NoMemory collapses under PO, while TextBuffer and StructMemory fully recover.}
\label{tab:cross_platform}
\begin{tabular}{l cc cc}
\toprule
& \multicolumn{2}{c}{FO} & \multicolumn{2}{c}{WeakPO} \\
\cmidrule(lr){2-3} \cmidrule(lr){4-5}
Memory & Success & Reward & Success & Reward \\
\midrule
NoMemory     & 100.0 & $-0.12$ & 15.6 & $-3.87$ \\
TextBuffer   & 100.0 & $-0.12$ & \textbf{100.0} & $-0.12$ \\
StructMemory & 100.0 & $-0.12$ & \textbf{100.0} & $-0.12$ \\
\bottomrule
\end{tabular}
\begin{flushleft}
\footnotesize FetchReach-v4 uses a 3-DOF goal-reaching task with the Fetch robot arm. The PO protocol and flash schedule (\texttt{flash\_every=5}, \texttt{flash\_len=1}) are identical to MetaWorld experiments. The identical NM$\to$TB/SM recovery pattern confirms that the layered diagnostic findings generalize beyond MetaWorld.
\end{flushleft}
\end{table}
